\scanpage{20}
\mspage{8}
% archive sheet number: 19
Il revient au même de dire que $\sigma_0,\sigma_1,\sigma_\infty,\rho$ sont des générateurs, avec relations de définition
\[
\overbrace{(\rho^{-1}\sigma_\infty\sigma_0\deleted{\rho})}^{l_\infty}
\overbrace{(\rho^{-1}\sigma_1\sigma_\infty\deleted{\rho})}^{l_1}
\overbrace{(\rho^{-1}\sigma_0\sigma_1\deleted{\rho})}^{l_0}=1.\tag*{(4)}
\]
\[
\sigma_0^2=\sigma_1^2=\sigma_\infty^2=1\tag*{(5)}
\]
\[
\rho^3=1.\tag*{(6)}
\]
% margin note written sideways; slanted
\margin{Il faut ajouter les relations $\rho\sigma_0\rho^{-1}=\sigma_1$ \unsure{etc.} On a deux gén. $\sigma_0,\rho$ en fait avec relations $\sigma_0^2=\rho^3=1$, la \deleted{$[\rho^{-1}\unsure{\sigma_0\rho\sigma_0\rho^{-1}}](\underbrace{\rho\sigma_0\rho^{-1}}_{\sigma_1})(\underbrace{\rho^{-1}\sigma_0\rho}_{\sigma_\infty})\unsure{\rho^{-1}}$} relation (4) étant conséquence des \unsure{précédentes} (\unsure{quand} on y fait $\sigma_1=\rho\sigma_0\rho^{-1}$, $\sigma_\infty=\rho^{-1}\sigma_0\rho$). Noter que $l_0=\deleted{\unsure{(\sigma_0\rho)^2}}\rho^{-1}(\sigma_0\rho)^2\rho=\boxed{\rho^{-1}\sigma_0\rho\sigma_0\rho^{-1}}$, $l_1=(\deleted{\rho}\sigma_0\rho)^2$, $l_\infty=(\rho^{-1}\sigma_0\deleted{\rho})^2$; $\operatorname{int}(\sigma_0)l_1=l_\infty$, $\operatorname{int}(\sigma_0)l_\infty=l_1$, $\operatorname{int}(\sigma_0)(l_0)=l_\infty^{-1}l_1^{-1}$.}
% slanted
\margin{\uline{NB} $\sigma_0(l_0)=l_\infty^{-1}l_1^{-1}$, $\sigma_0(l_1)=l_\infty$, $\sigma_0(l_\infty)=l_1\deleted{\unsure{l_\infty^{-1}}}$ et idem pour $\sigma_1,\sigma_\infty$.}

\uline{NB} \unsure{Si on} y fait \inserted{\uline{de plus}} $\sigma_0,\sigma_1,\sigma_\infty=1$, on trouve comme relation de définition $\rho^3=1$, \deleted{\unsure{alors}} sans plus, ce qui prouve que $\rho$ n'est pas un mot \unsure{p.v.|p.r.} aux $\sigma_i$\unsure{...}

Tout ceci est compatible avec les actions des lacets, \deleted{\unsure{et ces dernières ne seraient pas pour des|et ces dernières ne seraient pas des} \unsure{raisons}, et qu'on interprète $l'_i=l_i^{-1}$, $\sigma'_i=\sigma_{\unsure{\dot\rho|\rho}(i)}$} \deleted{\unsure{puis} $\sigma_i$} \deleted{\unsure{de même que} $\sigma_i'$} si on y remplace les $l_i$ par les $l_i^{-1}$. \unsure{Ici|Déjà} on a, comme $\sigma_i=\sigma'_{\dot\rho^{-1}(i)}$
\[
\left\{
\begin{aligned}
&\operatorname{int}(\sigma_i):\ \text{échange }l_j\text{ et }l_{\deleted{\unsure{k}}h}\quad \inserted{(i,j,h)\text{ distincts}}\text{ et}\\
&\qquad\deleted{\text{(et le dernier, }l_i,\text{ s'en déduit par (1))}}\\
&\operatorname{int}(\rho)l_i=l_{\dot\rho(i)}.
\end{aligned}
\right.\tag*{(7)}
\]

On voudrait maintenant expliciter les $\unsure{\varepsilon_i|\epsilon_i}$ avec $\varepsilon_i^2=l_i$ $(i=0,1,\infty)$, $\varepsilon_i$ étant un générateur du centralisateur (ou normalisateur, c'est \unsure{idem}) des $L_i$ dans $\unsure{\mathscr E^+|\mathcal E^+}$. (Ou le centralisateur dans $\mathscr E$). Les calculs faits plus haut, avec les chgts de notation indiqués, donnent
\[
\left\{\begin{aligned}
\varepsilon_0&=l_0^{\unsure{4|1|-1}}\sigma_1l_{\unsure{1|0|\infty}}^{\unsure{4|1|-1}}\\
\varepsilon_1&=l_\infty^{-1}\sigma_\infty\deleted{l_1^{-1}}l_1\\
\varepsilon_\infty&=l_0^{-1}\sigma_0l_\infty
\end{aligned}\right.\tag*{\unsure{(8)}}
\]
